% !TeX root = ../tesi.tex

\capitolononnumerato{Glossario}

Raccolgo in ordine alfabetico i termini tecnici introdotti nel testo. Alla prima occorrenza nel
corpo della relazione, il termine e il pedice G rimandano alla voce corrispondente.

\begin{description}
\voceglossario{adr}{ADR}{
        Documento che registra una decisione di architettura: il contesto in cui nasce, le
	alternative considerate, la scelta e le sue conseguenze. Conserva le motivazioni della scelta anche
	quando cambia chi lavora al progetto.
}

\voceglossario{api}{API}{
        Interfaccia con cui un programma rende disponibili dati o operazioni ad altri programmi,
	secondo un contratto definito. In Shopify, la Admin API permette di amministrare il negozio; la
	Storefront API espone le funzioni rivolte all’esperienza d’acquisto.
}

\voceglossario{b2b}{B2B}{
        Modello commerciale in cui un’azienda vende prodotti o servizi ad altre aziende. Nel
	commercio elettronico può richiedere cataloghi, listini e condizioni di acquisto diversi per
	ciascun cliente.
}

\voceglossario{b2c}{B2C}{
        Modello commerciale in cui un’azienda vende prodotti o servizi al consumatore finale. È il
	contesto dell’interfaccia conversazionale di ChatCart, rivolta a chi acquista nel negozio.
}

\voceglossario{crm}{CRM}{
        Insieme di pratiche e sistemi per gestire le relazioni di un’azienda con i propri clienti.
	Un sistema CRM raccoglie, per esempio, contatti, interazioni commerciali e richieste di assistenza.
}

\voceglossario{design-system}{\emph{Design system}}{
        Insieme coerente di principi, regole visive e componenti riusabili per progettare
	un’interfaccia. Comprende elementi come colori, tipografia, spaziature e comportamenti dei
	componenti, così da mantenere riconoscibile l’identità del prodotto o del marchio.
}

\voceglossario{e-commerce}{\emph{E-commerce}}{
        Commercio di prodotti o servizi attraverso canali digitali. Un negozio elettronico comprende
	il catalogo, la selezione dei prodotti, il carrello e il pagamento, oltre ai collegamenti con i
	sistemi che gestiscono ordini e logistica.
}

\voceglossario{erp}{ERP}{
        Sistema gestionale che integra dati e processi di diverse aree aziendali, come
	amministrazione, acquisti, magazzino e vendite. Il collegamento con il negozio permette, per
	esempio, di sincronizzare ordini e disponibilità dei prodotti.
}

\voceglossario{founder-senior-consultant}{\emph{Founder \& Senior Consultant}}{
        Ruolo che unisce la fondazione dell’azienda all’attività di consulenza esperta. In Halue lo
	ricopre Valentino Baraldo.
}

\voceglossario{graphql}{GraphQL}{
        Linguaggio di interrogazione per API e insieme di regole per eseguire le richieste su uno
	schema di dati tipizzato. Il programma che invia la richiesta specifica i campi di cui ha bisogno;
	Shopify lo adotta per la Admin API e la Storefront API.
}

\voceglossario{headless}{\emph{Headless}}{
        Architettura che separa la gestione dei dati e delle funzioni applicative dall’interfaccia
	che li presenta. In un negozio di questo tipo l’interfaccia è un’applicazione autonoma che comunica
	con la piattaforma commerciale tramite API; in un CMS la stessa separazione riguarda i contenuti e
	la loro presentazione.
}

\voceglossario{hydrogen}{Hydrogen}{
        \emph{Framework} di Shopify basato su React per costruire interfacce di negozi
	\emph{headless}. Offre strumenti e componenti per accedere ai dati del negozio tramite la
	Storefront API; Oxygen è il servizio di Shopify che può ospitare queste applicazioni.
}

\voceglossario{ipaas}{iPaaS}{
        Piattaforma erogata come servizio per collegare applicazioni e sistemi diversi. Offre
	connettori e strumenti per trasformare e trasferire i dati, riducendo la necessità di costruire
	un’integrazione distinta per ogni coppia di sistemi.
}

\voceglossario{junior-developer}{\emph{Junior Developer}}{
        Sviluppatore all’inizio del proprio percorso professionale, che realizza e verifica
	funzionalità con il supporto di colleghi più esperti.
}

\voceglossario{liquid}{Liquid}{
        Linguaggio di \emph{template} usato nei temi Shopify. Permette di combinare la struttura
	delle pagine con i dati del negozio, attraverso variabili, condizioni e cicli; la piattaforma
	elabora i modelli e genera le pagine da inviare al navigatore.
}

\voceglossario{llm}{LLM}{
        Modello linguistico di grandi dimensioni, addestrato su ampie raccolte di dati per elaborare
	e generare testo. In un agente può selezionare strumenti esterni di cui l’applicazione esegue le
	operazioni; il collegamento a questi strumenti gli permette di usare dati aggiornati che non fanno
	parte del suo addestramento.
}

\voceglossario{mcp}{MCP}{
        Protocollo aperto che standardizza la comunicazione tra applicazioni basate su modelli
	linguistici e servizi esterni. Un \emph{server} MCP può esporre strumenti, risorse e istruzioni; è
	l’applicazione che ospita il modello a gestire il collegamento e l’esecuzione delle richieste.
}

\voceglossario{project-management-assistant}{\emph{Project Management Assistant}}{
        Figura che supporta la gestione dei progetti, contribuendo all’organizzazione delle
	attività, all’aggiornamento della documentazione e al monitoraggio delle scadenze.
}

\voceglossario{project-service-manager}{\emph{Project/service manager}}{
        Responsabilità operativa di gestione di un progetto o di un servizio continuativo. Nel
	metodo di Halue cura la pianificazione delle attività e la capacità disponibile del gruppo di
	lavoro.
}

\voceglossario{pull-request}{\emph{Pull request}}{
        Proposta di integrare un insieme di modifiche da un \emph{branch} a un altro in un sistema
	di versionamento. Raccoglie le differenze nel codice, la discussione di revisione e gli esiti dei
	controlli automatici prima dell’integrazione.
}

\voceglossario{sal}{SAL}{
        Stato di avanzamento dei lavori. Nel contesto descritto indica un rapporto periodico che
	riassume il progresso del progetto, le scadenze, i rischi e i punti aperti, per condividere la
	situazione con le persone coinvolte.
}

\voceglossario{scrum}{Scrum}{
        \emph{Framework} per affrontare problemi complessi attraverso cicli brevi di lavoro,
	chiamati \emph{sprint}, e un confronto regolare sui risultati. Prevede responsabilità, eventi e
	artefatti definiti; nella \secref{sec:organizzazione-halue} distinguo questo riferimento
	dall’adattamento che Halue ne fa per il lavoro su commessa.
}

\voceglossario{semver}{SemVer}{
        Convenzione di versionamento semantico che esprime una versione nella forma
	\texttt{MAJOR.MINOR.PATCH}. Rispetto all’API pubblica dichiarata, il primo numero cambia per
	modifiche incompatibili, il secondo per nuove funzionalità compatibili e il terzo per correzioni
	compatibili.
}

\voceglossario{senior-project-manager}{\emph{Senior Project Manager}}{
        Figura con esperienza nella gestione dei progetti, responsabile del coordinamento di
	pianificazione, risorse e avanzamento delle attività.
}

\voceglossario{senior-software-engineer}{\emph{Senior Software Engineer}}{
        Professionista dello sviluppo del \emph{software} con esperienza consolidata, che affronta
	problemi progettuali complessi e contribuisce alle scelte tecniche e alla revisione del codice.
}

\voceglossario{skill}{\emph{Skill}}{
        Insieme riusabile di istruzioni e, quando necessari, risorse o strumenti, che guida un
	assistente basato su IA in un’attività specifica. Serve a rendere esplicito un procedimento, ma non
	garantisce da solo la correttezza dell’esecuzione.
}

\voceglossario{societa-benefit}{Società Benefit}{
        Società che affianca allo scopo di lucro una o più finalità di beneficio comune indicate
	nello statuto. La disciplina italiana richiede una gestione che bilanci questi obiettivi e una
	relazione annuale sull’impatto prodotto; la qualifica non coincide con una certificazione privata
	di sostenibilità.
}

\voceglossario{software-engineer}{\emph{Software Engineer}}{
        Professionista che applica metodi di ingegneria del \emph{software} alla progettazione,
	realizzazione, verifica e manutenzione di sistemi \emph{software}.
}

\voceglossario{solution-technical-architect}{\emph{Solution \& Technical Architect}}{
        Figura che traduce le esigenze applicative in una soluzione tecnica, definendo
	l’architettura del sistema, le tecnologie e le integrazioni tra componenti.
}

\voceglossario{sprint}{\emph{Sprint}}{
        Ciclo di lavoro di durata fissa, non superiore a un mese nel riferimento Scrum, durante il
	quale il gruppo persegue un obiettivo concordato. La revisione dei risultati e la retrospettiva
	permettono di adattare il prodotto e il modo di lavorare al ciclo successivo.
}

\voceglossario{staging}{\emph{Staging}}{
        Ambiente separato dalla produzione in cui collaudare una versione dell’applicazione prima di
	renderla disponibile agli utenti. Riproduce le condizioni di esercizio per quanto necessario alla
	verifica, senza coincidere con il servizio effettivamente usato dai clienti.
}

\voceglossario{storefront}{\emph{Storefront}}{
        Interfaccia di un negozio digitale rivolta al cliente, attraverso cui consultare il
	catalogo, selezionare prodotti e avviare l’acquisto. In Shopify può basarsi sui temi della
	piattaforma oppure su un’applicazione separata con architettura \emph{headless}.
}

\voceglossario{story-point}{\emph{Story point}}{
        Unità di stima relativa che un gruppo usa per confrontare le attività in termini di
	complessità, impegno e incertezza. Non corrisponde a un numero fisso di ore e acquista significato
	rispetto alle stime dello stesso gruppo; la scala di Fibonacci è una possibile convenzione, non un
	obbligo di Scrum.
}

\voceglossario{team-lead}{\emph{Team lead}}{
        Responsabilità di guida tecnica di un gruppo di lavoro. In Halue comprende l’analisi
	funzionale, le stime per il cliente e il coordinamento tecnico degli sviluppatori.
}

\voceglossario{ticket}{\emph{Ticket}}{
        Unità di lavoro registrata in un sistema di tracciamento, con un identificativo, una
	descrizione e uno stato. Può rappresentare una funzionalità, una modifica o un difetto, e
	raccogliere stime, discussioni e collegamenti alle modifiche del codice.
}

\voceglossario{ucp}{UCP}{
        Protocollo aperto per rendere interoperabili i servizi commerciali e gli agenti che operano
	per conto dell’utente. Definisce modalità condivise per scoprire le capacità di un negozio e
	interagire con le fasi dell’acquisto; riguarda il dominio commerciale, mentre MCP offre un
	meccanismo generale per collegare applicazioni e strumenti.
}

\voceglossario{ux-ui-digital-consultant}{\emph{UX/UI Digital Consultant}}{
        Consulente che progetta l’esperienza utente e l’interfaccia di un prodotto digitale,
	considerando percorsi d’uso, struttura delle informazioni e aspetto visivo.
}
\end{description}
